\section{議論}\label{sec:discussion}

\subsection{何が確立し、何が仮説か}
本論文で確立したのは次の 3 点である。
\begin{enumerate}
\item refined クロスキャップ不変量 $g = f - \tfrac12 f^{\cov}$ は、一様な規格化と $M(q^{-1})$ の hatted のもとで、調べた範囲（unknot、Hopf 絡み目、トーラス結び目の表現 [1]、trefoil の 2 箱・3 箱、$T(2,5)$ の 2 箱）のすべてで整数係数の Laurent 多項式になり、$t = q$ で Mariño の結果に戻る【計算】。
\item 全 $d$ の構造 $g = \tfrac12 E/(1+t) + C$ と、$E$ が偶数のセクターにしか現れないこと、$E$ の境界因子 $z_q^{n-1}(z_q^+)^n$ は、仮説 \ref{hyp:classes} のもとで導出され【導出】、計算と一致する【計算】。
\item 被覆の brane--brane セクターは \cite{KN17} の不変量そのもので、brane--image セクターは状態の空間として同定できる【計算】。
\end{enumerate}
仮説として残るのは、被覆の部分が A 型の refined 不変量で与えられること（仮説 \ref{hyp:hg}。被覆を引いた後の $g$ の整数性と、極の打ち消しの一意性で裏付けた）、一様な規格化が framing 0 であること、一粒子状態の分類（仮説 \ref{hyp:classes}）である。
予想として掲げるのは、表現 [1] の判定式（予想 \ref{conj:criterion}、$m \le 5$ で確認）と、普遍因子の構造（予想 \ref{conj:structure}）である。

\subsection{文献 \texorpdfstring{\cite{KN17}}{[KN17]} との関係}\label{sec:kn}
ここでの正値性は §\ref{sec:criterion} の判定式の意味、すなわち $G$ 展開の後に $\lambda^\beta$ のブロックごとに符号がそろうことである（§\ref{sec:positivity} の冒頭）。
被覆（向き付けられた部分）の brane--brane のセクターは \cite{KN17} の U$(N_A)$ の refined LMOV 不変量に全係数で一致し（§\ref{sec:split}）、したがって \cite{KN17} の $G$ 展開と正値性をそのまま満たす（\ref{app:cover}）。
brane--image のセクターは、境界因子を外すと切り替わりのブロック $k = 2m$ の $-S\otimes S$ を除いてそろい、その混在は対の状態の引き算として理解できる（§\ref{sec:split}）。
破れが残るのはクロスキャップ $g$ の多箱で、これは同じ普遍因子 $(\bt^2 - 1)$ をもつ $\tilde I$ の固定点のセクターである。
根拠は trefoil と $T(2,5)$ の 2 箱まで（brane--image は $T(2,7)$ まで）で、3 箱は brane と image に分けていない。
したがって、refined の Kauffman 不変量の正値性の問題は、\cite{KN17} の枠組みに orientifold の偶奇の構造（$\tilde I$ の符号と、凍結した状態の多重項の基底）を加える問題に帰着する（予想）。

\subsection{未解決の問題}
\begin{enumerate}
\item \textbf{凍結した状態の多重項}：多箱の $c = 1$ 部分は $\sigma$ 対称でなく、$G$ では展開できない。$G = T_+T_-$ の因子 $T_\pm$ を使っても、$\sigma$ 対称でない部分は展開できるようにならない（§\ref{sec:index} の注意）。
計算結果 \ref{cr:multibox-sigma} により、非対称性のうち unrefined の部分は $c$ の偶奇の規則で取り除け、残る refined の $\sigma$ 奇部分は $(\bt^2-1)(\ql - \ql^{-1})$ で割り切れる。$\sigma$ 偶・奇の分解は半整数で、$T(2,5)$ では符号もそろわないので、この分解は正値性の判定には足りない。
$G$ 展開は向き付けられた曲面の構造（$H^1$ の Hodge 分解に働く Lefschetz の $\mathfrak{sl}_2$）から来るので、向き付けられない曲面に巻く凍結した状態にそれを期待する根拠は無い（§\ref{sec:index}）。向き付けられない面の refined のハンドルとクロスキャップの因子を幾何から決めることが、第一の課題である。
\item \textbf{$E$ と $C$ の分離}：指数からは決まらない（命題 \ref{prop:EC}）。Hopf の $(4,4)$ と結び目の 2 箱で未決定である（brane 2・image 2 の trefoil では $\lambda^{-12}$ まで決まり、$\lambda^{-14}$ 以降が未決定。\ref{app:n2}）。
境界ゼロモードへの作用は決まった（命題 \ref{prop:boundary}）。残るのはモジュライ空間の部分で、$\tilde I = I_X\otimes e^{i\pi J_2}$ の符号を決めるには状態の $t$ 電荷を $J_2$ と $S_R$ に分ける必要があり、指数からは得られない。\cite{Young16} の cohomological Hall module が CoHA 上で自由でないことは、この不定性の代数的な形である（loop quiver では決まり、framed unknot で $E$ が一意になることと整合する）。
\item \textbf{普遍因子の意味}：$(\ba^2 - \bq^2)$ は SO(2) の点で、$(\bt^2 - 1)$ は $t = q$ と偶奇で説明できる。ただし、$\lambda \to \tlam$ のずれの向きが一意に決まる理由と、閉弦の Kähler 変数の三重化との関係は未解明である。
\item \textbf{他の結び目}：非トーラス結び目では、refined SO の値に定義が無い（refined CS の $S, T$ が使えない）。
固定 $N$ では、SO(4) $=$ SU(2)$^2$、SO(6) $=$ SU(4) の同一視で A 型の入力（\cite{KNTZ19} の長方形彩色の式など）を使える。二重ツイスト結び目 7 個で、$N = 2$ の SO 側（$|R| \le 3$）の hatted の整数性と一様な規格化が成り立つことを確かめた。
大きな $N$ の値には、D 型の DAHA の iterated torus 結び目への拡張が要る。
\item \textbf{B 型と高い彩色の Sp}：Sp$(2N)$ は §\ref{sec:sp} で扱った（表現 [1] と trefoil の 2 箱）。$T(2,5)$ の 2 箱は補間に $N \le 13$ が要り、現在の実装では計算量が大きすぎる。閉弦の B 型（\cite{AM22}）は $c = 1$ だけで、分母に $(t^{n/2} + t^{-n/2})$ をもつ。命題 \ref{prop:alld} の分類で B 型をどう読むかは調べていない。
\item \textbf{Hilbert scheme の模型}：\cite{ORS12} のセル分解（$\Gamma = \langle 2, 3\rangle$）を実装すると、$T(2,3)$ の HOMFLY ホモロジー（\cite{DGR05} の unreduced superpolynomial）が $q^{2l}$ の 16 次まで再現され、$(q, t, a)$ の次数とセルの（長さ、重み、生成元の数）の辞書が固定できる。
一方、固定セクターの指数 $P = (1-t)x$ は $-(\ba^2 - 1)(\ba^2 - \bq^2)(\bq^2\bt^2 - 1)(\ba^2\bq^2\bt^2 - 1) N_8/(\ba^{12}\bq^4\bt^{10})$ と因数分解し（$N_8$ は 8 項、framed unknot と同じ境界因子 $(\lambda^2-1)(q-1)$ をもつ）、被覆の表現 [1] の $\psi^2$ では書けない。$E$ を幾何から決めるには、brane--image の composite セクターのセル模型と、その上の実構造が要る（計算と予想）。
\end{enumerate}

\subsection{今後}
本論文の定義と構造は、トーラス結び目の表現 [1] では閉じた形の判定式まで到達し、多箱では整数性と普遍因子まで到達した。
次の段階は、凍結した状態の多重項の構造を向き付けられない曲面の幾何から決めることと、D 型の DAHA を iterated torus 結び目に広げて定義の範囲を増やすことである。
前者が決まれば、§\ref{sec:positivity} の多箱の符号の問題は、判定式の多箱版として定式化できる。
